一刚性轴对称细长卫星沿半径 \(R\) 的牛顿引力圆轨道运行，中心引力参数 \(\mu\)，轨道平均角速率 \(n^2=\mu/R^3\)。卫星尺度远小于 \(R\)，忽略轨道反作用；对称轴在轨道平面内，与当地径向夹角为 \(\theta\)，绕轨道法向的转动惯量为 \(I_\perp>0\)，沿对称轴为 \(I_\parallel\)，取 \(I_\parallel<I_\perp\)，无外控转矩。仅保留引力势的四极项。
(1) 从中心引力泰勒展开推导姿态相关的重力梯度势能，说明沿径向与沿切向放置的能量差。
(2) 在随轨道旋转的坐标下写姿态拉格朗日量，求非线性俯仰方程及径向定向的小扰动频率。
(3) 构造旋转系Jacobi型守恒积分与分界能量，求从径向平衡以纯角速度启动时越过切向不稳定点的临界速度。
(4) 检查球形惯量极限、轨道频率极限以及扭矩的符号；区分该无阻尼保守稳定与受姿态控制时的渐近稳定。
